\documentclass[../../../include/open-logic-section]{subfiles}

\begin{document}

\olfileid[tr]{sth}{cardinals}{hp}
\olsection{Ek: Hume İlkesi}

\olref[cp]{sec} içinde Cantor İlkesi'ni betimledik. İlke şuydu:
\begin{align*}
	\card{A} = \card{B} & \text{ ancak ve ancak } A \approx B.
\intertext{Bu, günümüzde \emph{Hume İlkesi} denen şu ilkeye çok benzer:}
	\fregenum{x} {F(x)} = \fregenum{x}{G(x)} & \text{ ancak ve ancak } F \sim G.
\end{align*}
Burada `$F\sim G$', tam olarak $F$ler kadar $G$ bulunduğunu, yani $F$lerle
$G$ler arasında bire bir eşleme kurulabildiğini kısaltır:
\begin{align*}
	\exists R(&\forall v\forall y(Rvy \lif (Fv \land Gy)) \land {}\\
		&\forall v(Fv \lif \lexists![y][Rvy]) \land {}\\
		&\forall y(Gy \lif \lexists![v][Rvy])).
\end{align*}
Fakat Hume İlkesi ile Cantor İlkesi arasında bir tür farkı vardır. Cantor
İlkesi'nin ifadesindeki ``$A$'' ve ``$B$'' değişkenleri \emph{kümeleri}
temsil eden birinci dereceden terimlerdir. Hume İlkesi'nin ifadesindeki
``$F$'', ``$G$'' ve ``$R$'' ise birinci dereceden terimler \emph{değildir};
bunun yerine \emph{yüklem konumunda} bulunurlar. (Belki de
\emph{özellikleri} temsil ederler.) Dolayısıyla Hume İlkesi'ni Türkçe şöyle
açıklayabiliriz: $F$lerin sayısı $G$lerin sayısına ancak ve ancak $F$lerle
$G$ler arasında bire bir eşleme varsa eşittir. Buna \emph{Hume İlkesi}
denir; çünkü Hume bir keresinde şunu yazmıştır:
\begin{quote}
  İki sayı, birinin her birimine ötekinin bir birimi karşılık gelecek
  biçimde birleştirildiğinde onların eşit olduğunu söyleriz.
  \citep[Kısım III, Kitap 1, \S1]{Hume1740}
\end{quote}
Hume İlkesi, Hume'dan bu parçayı alıntılayıp ardından (gerçekte) bizim Hume
İlkesi dediğimiz şeyin taslağını çizen \citet[\S63]{Frege1884} tarafından
çağdaş matematiksel-mantıksal tartışmada öne çıkarılmıştır.

Hume İlkesi ile Temel Yasa V arasındaki yapısal benzerliğe dikkat
etmelisiniz. Temel Yasa V'yi \olref[story][blv]{sec} içinde şöyle
formüle ettik:
\[
	\fregeext{x}{F(x)} = \fregeext{x}{G(x)} \text{ ancak ve ancak }
	\lforall[x][(F(x) \liff G(x))].
\]
Bu noktada biraz yorum ve karşılaştırma yararlı olabilir.

Hume İlkesi ya da Temel Yasa V gibi bir ilkeyi iki biçimde ele alabiliriz:
\emph{yüklemsel} ya da \emph{yüklemsel olmayan}
(\olref[story][predicative]{sec} kesimini hatırlayın). Temel Yasa V'nin
yüklemsel olmayan okumasında her $F$ için $\fregeext{x}{F(x)}$ nesnesi,
Temel Yasa V'yi formüle ederken kullandığımız niceleme evreninin içinde yer
alır. Benzer biçimde Hume İlkesi'nin yüklemsel olmayan okumasında her $F$
için $\fregenum{x}{F(x)}$ nesnesi, Hume İlkesi'ni formüle ederken
kullandığımız niceleme evreninin içinde yer alır. Buna karşılık
\emph{yüklemsel} anlayışta $\fregeext{x}{F(x)}$ ve
$\fregenum{x}{F(x)}$ nesneleri farklı bir evrenden gelen varlıklar olur.

Temel Yasa V'yi yüklemsel olmayan biçimde okursak Naif Küme Oluşturma
aracılığıyla tutarsızlığa yol açar (ayrıntılar için
\olref[story][blv]{sec} kesimine bakın). Naif Küme Oluşturma'ya çok benzer
biçimde, onu \emph{yüklemsel} okuyarak tutarlı hâle getirebiliriz. Fakat o
zaman muhtemelen istediğimiz her şeyi yapmayacaktır.

Hume İlkesi ise yüklemsel olmayan biçimde \emph{tutarlı olarak} okunabilir.
Bu biçimde okunduğunda oldukça güçlüdür.

Örneğin, açıkça hiçbir şeyin sağlamadığı ``$x\neq x$'' yüklemini ele alın.
Hume İlkesi bize şimdi bir $\#x(x\neq x)$ nesnesi verir. Bunu $0$ sayısı
olarak ele alabiliriz. \emph{Yüklemsel olmayan} anlayışta---ve
\emph{yalnızca} bu anlayışta---bu $0$ varlığı ilk niceleme evrenimizin
içinde bulunur. Dolayısıyla ``$x=0$'' yüklemiyle Hume İlkesi'ni anlamlı
biçimde uygulayıp bir $\#x(x=0)$ nesnesi elde edebiliriz. Bunu $1$ sayısı
olarak ele alabiliriz. Ayrıca Hume İlkesi $0\neq1$ olduğunu gerektirir;
çünkü kendisiyle özdeş olmayan nesnelerle $0$'la özdeş nesneler arasında
bire bir eşleme kurulamaz (birincilerden hiç yoktur, ikincilerden bir tane
vardır). Yine yüklemsel olmayan biçimde çalıştığımızda $1$ ilk niceleme
evrenimizde bulunur. Bu yüzden ``$(x=0\lor x=1)$'' yüklemiyle Hume
İlkesi'ni anlamlı biçimde uygulayıp $\#x(x=0\lor x=1)$ nesnesini elde
edebiliriz. Bunu $2$ sayısı olarak ele alabilir ve $0\neq2$, $1\neq2$ vb.
olduğunu gösterebiliriz.

Kısacası yüklemsel olmayan biçimde ele alınan Hume İlkesi \emph{sonsuz
sayıda nesne} bulunduğunu gerektirir. Bu, \emph{yeni-Fregeci
mantıkçıları} Hume İlkesi'ni aritmetiğin temeli olarak almaya teşvik
etmiştir.

Ancak Frege'nin \emph{kendisi} Hume İlkesi'ni aritmetiğin temeli olarak
almadı. Bunun yerine Hume İlkesi'ni açık bir tanımdan ispatladı:
$\fregenum{x}{F(x)}$, $F\sim\Phi$ kavramının kaplamı olarak tanımlanır.
Çağdaş terimlerle bunu
$\fregenum{x}{F(x)}=\Setabs{G}{F\sim G}$ olarak ifade etmeye çalışabiliriz;
fakat bu bizi Naif Küme Oluşturma'nın sorunlarına geri götürür.

\end{document}
